वस्तूंच्या संग्रहाला \emph{संच} म्हणतात आणि त्या वस्तू त्या संचाचे घटक असतात. $x$ हा $A$ चा घटक असेल, तर आपण $x \in A$ असे लिहितो. संच \emph{विस्तारात्मक} असतात, म्हणजेच ते त्यांच्या घटकांवरून पूर्णपणे निर्धारित होतात. घटकांची स्पष्ट \emph{यादी देऊन} किंवा त्यांचा समान गुणधर्म देऊन (\emph{गुणधर्माधारित संचरचना}) संच निर्दिष्ट करता येतात. विस्तारात्मकतेचा अर्थ असा की घटक कोणत्या क्रमाने किंवा कोणत्या पद्धतीने दिले आहेत याने फरक पडत नाही. $A$ आणि~$B$ हे एकच संच आहेत ($A = B$) हे सिद्ध करण्यासाठी $A$ चा प्रत्येक घटक $B$ चा घटक आहे आणि उलटही तसेच आहे हे सिद्ध करावे लागते.

महत्त्वाच्या संचांमध्ये नैसर्गिक संख्या ($\Nat$), पूर्णांक ($\Int$), परिमेय संख्या ($\Rat$) आणि वास्तव संख्या ($\Real$) यांचे संच येतात. वस्तूंच्या \emph{चिन्हमाला} ($X^*$) आणि अनंत \emph{अनुक्रम} ($X^\omega$) यांचे संचही महत्त्वाचे आहेत. $A$ चा प्रत्येक घटक $B$ चाही घटक असेल, तर $A$ हा $B$ चा \emph{उपसंच} आहे, म्हणजे $A \subseteq B$. संच~$B$ च्या सर्व उपसंचांचा संग्रह हाही एक संच आहे. तो $B$ चा \emph{घातसंच}~$\Pow{B}$ आहे. संचांचा \emph{संयोग} $A \cup B$ आणि \emph{छेद} $A \cap B$ घेता येतो. \emph{क्रमित जोडी}~$\tuple{x, y}$ मध्ये $x$ आणि~$y$ या दोन वस्तू नेमक्या याच क्रमाने येतात. $\tuple{x, y}$ आणि $\tuple{y, x}$ या भिन्न जोड्या आहेत (जर $x = y$ नसेल तर). $x \in A$ आणि $y \in B$ असलेल्या सर्व $\tuple{x, y}$ जोड्यांच्या संचाला $A$ आणि $B$ यांचा \emph{कार्तीय गुणाकार}~$A \times B$ म्हणतात. $A \times A$ साठी आपण $A^2$ असे लिहितो. उदाहरणार्थ, $\Nat^2$ हा नैसर्गिक संख्यांच्या जोड्यांचा संच आहे.
